\documentclass[10pt,a4paper]{amsart}
\usepackage{amsmath,amssymb,amsthm,graphicx}
% Retained minimal preamble for the Stacks Project chapter more-algebra.tex (More on Algebra).
% The chapter file pinned-original-source.tex carries only \input{preamble} on its first line, so no
% preamble text lives inside the chapter itself.  The Stacks Project uses one single shared preamble for
% every chapter; the snapshot stages it identically as sources/stacks-chow/preamble.tex and
% sources/stacks-etale-cohomology/preamble.tex.  The block below is copied mechanically, line for line,
% from sources/stacks-chow/preamble.tex lines 14-18 (the Xy-pic setup) and 154-217 (the native theorem
% declarations together with the mathematical macros).  Nothing was added, removed or redefined; the
% wrapper class and the amsmath/amssymb/amsthm/graphicx base packages are supplied by the assembler.
% Chapter: more-algebra.tex -- More on Algebra.
% For commutative diagrams we use Xy-pic
\usepackage[all]{xy}

% We use 2cell for 2-commutative diagrams.
\xyoption{2cell}

% Theorem environments.
%
\theoremstyle{plain}
\newtheorem{theorem}[subsection]{Theorem}
\newtheorem{proposition}[subsection]{Proposition}
\newtheorem{lemma}[subsection]{Lemma}

\theoremstyle{definition}
\newtheorem{definition}[subsection]{Definition}
\newtheorem{example}[subsection]{Example}
\newtheorem{exercise}[subsection]{Exercise}
\newtheorem{situation}[subsection]{Situation}

\theoremstyle{remark}
\newtheorem{remark}[subsection]{Remark}
\newtheorem{remarks}[subsection]{Remarks}

\numberwithin{equation}{subsection}

% Macros
%
\def\lim{\mathop{\mathrm{lim}}\nolimits}
\def\colim{\mathop{\mathrm{colim}}\nolimits}
\def\Spec{\mathop{\mathrm{Spec}}}
\def\Hom{\mathop{\mathrm{Hom}}\nolimits}
\def\Ext{\mathop{\mathrm{Ext}}\nolimits}
\def\SheafHom{\mathop{\mathcal{H}\!\mathit{om}}\nolimits}
\def\SheafExt{\mathop{\mathcal{E}\!\mathit{xt}}\nolimits}
\def\Sch{\mathit{Sch}}
\def\Mor{\mathop{\mathrm{Mor}}\nolimits}
\def\Ob{\mathop{\mathrm{Ob}}\nolimits}
\def\Sh{\mathop{\mathit{Sh}}\nolimits}
\def\NL{\mathop{N\!L}\nolimits}
\def\CH{\mathop{\mathrm{CH}}\nolimits}
\def\proetale{{pro\text{-}\acute{e}tale}}
\def\etale{{\acute{e}tale}}
\def\QCoh{\mathit{QCoh}}
\def\Ker{\mathop{\mathrm{Ker}}}
\def\Im{\mathop{\mathrm{Im}}}
\def\Coker{\mathop{\mathrm{Coker}}}
\def\Coim{\mathop{\mathrm{Coim}}}

% Boxtimes
%
\DeclareMathSymbol{\boxtimes}{\mathbin}{AMSa}{"02}

%
% Macros for moduli stacks/spaces
%
\def\QCohstack{\mathcal{QC}\!\mathit{oh}}
\def\Cohstack{\mathcal{C}\!\mathit{oh}}
\def\Spacesstack{\mathcal{S}\!\mathit{paces}}
\def\Quotfunctor{\mathrm{Quot}}
\def\Hilbfunctor{\mathrm{Hilb}}
\def\Curvesstack{\mathcal{C}\!\mathit{urves}}
\def\Polarizedstack{\mathcal{P}\!\mathit{olarized}}
\def\Complexesstack{\mathcal{C}\!\mathit{omplexes}}
% \Pic is the operator that assigns to X its picard group, usage \Pic(X)
% \Picardstack_{X/B} denotes the Picard stack of X over B
% \Picardfunctor_{X/B} denotes the Picard functor of X over B
\def\Pic{\mathop{\mathrm{Pic}}\nolimits}
\def\Picardstack{\mathcal{P}\!\mathit{ic}}
\def\Picardfunctor{\mathrm{Pic}}
\def\Deformationcategory{\mathcal{D}\!\mathit{ef}}

% --- appended 2026-10-12 (Chief): the project own definitions of the reference/slogan/history
% display environments, copied mechanically from the shared Stacks preamble lines 8-12
% (\usepackage{verbatim} plus the three \newenvironment lines, which the chapter preamble keeps
% verbatim from the project preamble). Nothing added, removed or redefined: these environments
% carry no mathematical content and the project itself typesets them as comments. Without these
% three lines a counted statement whose own text uses them fails with Environment slogan undefined.
\usepackage{verbatim}
\newenvironment{reference}{\comment}{\endcomment}
\newenvironment{slogan}{\comment}{\endcomment}
\newenvironment{history}{\comment}{\endcomment}
\usepackage{cleveref}
\crefname{originallemma}{original lemma}{original lemmas}
\Crefname{originallemma}{Original lemma}{Original lemmas}
\makeatletter
\newlabel{derived-lemma-two-ss-complex-functor}{{015J}{source}{}{sourcepointer-1}{}}
\newlabel{derived-lemma-two-ss-complex-functor@cref}{{[originallemma][0][]015J}{[1][0][]source}{}{}{}}
\makeatother
\title{For a four-term inverse system of complexes the zeroth cohomology of the limit is the limit of the zeroth cohomologies under two vanishing hypotheses}
\author{The Stacks Project Authors}
\date{stacks-project:more-algebra.tex:local-snapshot-20261012; GNU FDL 1.2 or later}
\begin{document}
\maketitle
\setcounter{subsection}{1}
\noindent\textit{Editorial note.} All delivered mathematical text is the author's own, copied from the pinned Stacks Project chapter file more-algebra.tex; no statement or proof step was rewritten, reordered or removed.\par\medskip
\noindent\textit{Editorial note.} The chapter file pinned-original-source.tex carries only \textbackslash{}input\{preamble\} on its first line, so no preamble text lives inside the chapter. The retained file sources/stacks-more-algebra/retained-preamble.tex is a mechanical, line-for-line copy of the project's single shared preamble as staged in the snapshot; nothing was added, removed or redefined and no package was added.\par\medskip
\noindent\textit{Editorial note.} The source URL is the source-level Stacks Project URL, as required when more than one card is produced from one source; the per-claim tag is carried by claim\_key.\par\medskip
\noindent\textit{Editorial note.} The counted statement span contains the author's four-term display, delivered verbatim. The counted proof span contains the author's footnote on injectivity, including its 'Details omitted.', and two displayed pages written as matrices; all are delivered verbatim. The author writes 'A inverse system' in the footnote, reproduced unchanged. The single reference inside the counted proof span is recorded as a non-load-bearing navigation pointer.\par\medskip
\noindent\textit{Allowed standard prerequisites (editorial).}\begin{itemize}
\item The right derived functor R lim on D(Ab(N)) and the first derived limit R\textasciicircum{}1 lim. used throughout the counted proof span. Supplies the terms appearing in both spectral sequences.
\item The two spectral sequences computing the cohomology of a complex of complexes under a functor. invoked at proof lines 23044-23050. Are the two computations whose comparison yields the conclusion.
\item Injectivity in the category of inverse systems of abelian groups. explained in the footnote at proof lines 23045-23049. Legitimises the use of the spectral sequences by providing enough injectives.
\end{itemize}\medskip
\ifdefined\hypertarget\hypertarget{sourcepointer-1}{}\fi
\noindent\textit{Source reference.} Original lemma 015J: focused-seeds/AGP/sources/stacks-snapshot-20261011/tags.txt, source lines 1510--1510. Cross-chapter navigation pointer into the Stacks Project; the standard fact it supplies is applied inside the counted proof span and the argument resting on it stays there.\par\medskip
\begin{lemma}
\label{lemma-apply-Mittag-Leffler-again}
Let
$$
(A^{-2}_n \to A^{-1}_n \to A^0_n \to A^1_n)
$$
be an inverse system of complexes of abelian groups and denote
$A^{-2} \to A^{-1} \to A^0 \to A^1$ its limit. Denote
$(H_n^{-1})$, $(H_n^0)$ the inverse systems of cohomologies, and
denote $H^{-1}$, $H^0$ the cohomologies of $A^{-2} \to A^{-1} \to A^0 \to A^1$.
If
\begin{enumerate}
\item $(A^{-2}_n)$ and $(A^{-1}_n)$ have vanishing $R^1\lim$,
\item $(H^{-1}_n)$ has vanishing $R^1\lim$,
\end{enumerate}
then $H^0 = \lim H_n^0$.
\end{lemma}

\begin{proof}
Let $K \in D(\textit{Ab}(\mathbf{N}))$ be the object represented
by the system of complexes whose $n$th constituent
is the complex $A^{-2}_n \to A^{-1}_n \to A^0_n \to A^1_n$.
We will compute $H^0(R\lim K)$ using both spectral
sequences\footnote{To use these spectral sequences we have to
show that $\textit{Ab}(\mathbf{N})$ has enough injectives.
A inverse system $(I_n)$ of abelian groups is injective if and only
if each $I_n$ is an injective abelian group and the transition maps are
split surjections. Every system embeds in one of these. Details omitted.} of
Derived Categories, Lemma \ref{derived-lemma-two-ss-complex-functor}.
The first has $E_1$-page
$$
\begin{matrix}
0 & 0 & R^1\lim A^0_n  & R^1\lim A^1_n \\
A^{-2} & A^{-1} & A^0  & A^1
\end{matrix}
$$
with horizontal differentials and all higher differentials are zero.
The second has $E_2$ page
$$
\begin{matrix}
R^1\lim H^{-2}_n & 0 & R^1\lim H^0_n & R^1 \lim H^1_n \\
\lim H^{-2}_n &
\lim H^{-1}_n &
\lim H^0_n &
\lim H^1_n
\end{matrix}
$$
and degenerates at this point. The result follows.
\end{proof}

\end{document}
